\subsection{Additional Cross-Dataset Component Validation}\label{sec:gpuvalidation}
The additional experiments use the same four datasets and fixed splits,
seeds 42, 123 and 777, the largest candidate width (64, or 256 for
CIFAR-10), and 300 epochs. There are 12 ARM runs and 24 ARD runs,
including two prior-centering settings, followed by 12 native-pruning
evaluations of the sample-centred ARD variant. All use the same RTX 3070,
Python, PyTorch, CUDA, cuDNN, driver and operating-system environment
listed in Section~\ref{sec:provenance}. They are separate from the
five-seed historical comparison and the one-seed MNIST capacity study.
The three-seed summaries use sample SDs; the repeated seeds and shared
splits do not constitute independent dataset replications.

The ARM implementation uses paired Bernoulli gates and per-example
loss differences before averaging the gate-gradient estimate. A supplied
linear-loss test checks this estimator against its analytic gradient.
The reconstruction constraint uses summed squared error for every
dataset, including dSprites. Its threshold is the ordinary $\beta=1$
anchor's deterministic validation MSE multiplied by 1.10 and by the
number of pixels. The network loss includes a gate-masked KL sum and
a constraint multiplier, and the expected gate-count penalty is applied
when the moving-average constraint is satisfied. Native evaluation
uses a binary mask with gate probability greater than 0.5. We report
posterior-mean validation/test MSE and validation MSE averaged over
32 posterior draws, both with that mask.

This is component validation, not exact reproduction of the source
algorithm~\cite{deboom2020geco}. The local gates start at probability
0.5 and update immediately, rather than remaining open until the first
constraint hit. The local KL is a masked sum rather than the source's
active-count-normalised quantity. The anchor-derived threshold and
local gate-update settings also differ. A failed local run therefore
does not establish a failure of the published method. Endpoint training
constraint satisfaction and deterministic validation $Q$ are different
events and are tabulated separately.

The ARD implementation fixes its training weight at one and reserves
the first 1800 of 18,000 training-pool examples for prior updates; the
remaining 16,200 supply parameter-gradient updates. The prior-update
pool is fixed, with latent samples refreshed at epochs 1, 11, \ldots,
291 and after training. This differs from the source's resampled pool
and tuned regularisation weight~\cite{saha2025ardvae}. We evaluate both
a zero centre in the variance update, as set in the source, and a
sample-centred variance-update variant; each is trained separately.
Both retain zero prior mean in the Gaussian closed-form KL.

At the final model, automatic differentiation evaluates the decoder
Jacobian at posterior means for the first 100 validation examples.
For prior variance $v_j$ and decoder output coordinates $\ell$, the score is
\begin{equation}
 w_j=\frac{1}{100}\sum_{i=1}^{100}\sum_{\ell=1}^{p}
 \left(\frac{\partial g_\ell}{\partial z_j}(\mu(x_i))\right)^2,
 \qquad s_j=w_jv_j.
 \label{eq:exactard}
\end{equation}
The readout is the number of descending scores required to explain
99\% of their total. Exact refers to automatic differentiation instead
of finite perturbations; it does not remove sample or training uncertainty.
The initial ARD comparison evaluates the unpruned full-width model.
The subsequent masking experiment below evaluates the retained axes.
For both methods,
dimension transfer maps the count to the nearest candidate width and
evaluates the cached ordinary $\beta=1$ VAE with the same seed.

The initial 12 ARM and 24 ARD runs retain evaluation-mode MSE curves every five epochs.
ARM uses a 2000-example training subset; the ARD training-pool subset
includes its reserved prior-update examples and is identified as such.
Recorded training time includes these periodic evaluations. ARM also
records final native-evaluation time; ARD separately records Jacobian
time but not final native-evaluation time. These are partial timings.
The additional scripts do not save model checkpoints.

\paragraph{Native ARD masking and fixed-budget decoder adaptation.}
The supplied internal specification fixes four conditions before the
new masking runs: full model, nonselected-axis masking, decoder-only
adaptation, and transfer. It is not a public preregistration.
The training script repeats the sample-centred variance-update setting
with the same three seeds per dataset. Its full-model validation/test
MSE and selected counts match the earlier sample-centred records exactly;
these are paired extensions of those settings, not independent baseline
replications. For each run the highest-scoring axes accounting for
99\% of total relevance are retained and
all other coordinates are set to the zero prior mean. Masking preserves
the network dimensions; it does not establish a smaller compiled model
or an inference-speed gain. The mask is computed once and is not
reselected after adaptation.

The third condition starts from that same masked model and updates only
decoder parameters for 30 epochs, with Adam, learning rate $10^{-3}$ and
batch size 128. Encoder parameters and the prior remain fixed.
The objective is the original reconstruction term on masked posterior
samples. Adaptation uses all 18,000 training-pool examples, including
the 1800 reserved for prior updates during initial training; this is
additional data exposure and computation relative to the masked
condition. No validation early stopping is applied. All four primary
conditions use posterior-mean validation and test MSE; the mask-only
condition also has a 32-draw validation-MSE diagnostic. The nearest-grid
transfer uses existing ordinary-VAE cache records, not a further model
trained by the masking script. Test metrics do not determine the mask
or the fixed adaptation budget.

The Jacobian score is still Equation~\eqref{eq:exactard}. To bound memory,
CIFAR-10 uses batches of four examples and dSprites accumulates squared
output-coordinate gradients in batches of eight instead of materialising
the full Jacobian. This changes the computational route, not the score
definition. The new JSON retains counts and MSE but not axis identities,
score vectors, training curves or model checkpoints. Timing and
record-level limits are detailed in Appendix~\ref{app:gpu}.

An evaluation-code audit limits use of the additional ELBO fields.
The 32-sample option changes MSE only: reconstruction likelihood is still
evaluated at the posterior mean, while KL uses an unmasked standard-normal
prior, including for ARD. These fields are not Monte Carlo ELBOs for the
trained native models. The native-masking script inherits this issue,
despite its specification requesting true ELBO. We exclude all such
fields and use the additional records for MSE, gates, relevance counts
and available curves; the capacity-selection ELBO evidence
remains the separate MNIST experiment in Section~\ref{sec:newvalidation}.
